% Budget: 0.20 pp (AGENT) / 0.15 pp (ICRITO).
% Future work: at most two sentences (Kill-and-Close framework is paper 2).
\section{Conclusion}\label{sec:conclusion}
Containing an agent ends its future actions; closing the incident requires accounting for its past ones. A runbook that equates a sent cancel with a stopped payment reports closures the processor contradicts. Requiring processor-side evidence removes those false closures, and in this model the price of that evidence depends on when it is sought: waiting for it before each next cancel costs loss under scarce capacity, collecting it after all cancels are sent costs little. Remediation that closes only on evidence also resumes safely after a crash, given a journal and idempotent cancels. Code and trial data are available at \ArtifactURL, and the per-run data behind every verdict at \DataLink. Joining revocation and disposition into one verified incident procedure is the next step.
